\documentclass[12pt]{article}
\usepackage[T1]{fontenc}
\usepackage[utf8]{inputenc}
\usepackage{lmodern}
\usepackage{textcomp}
\usepackage{enumitem}
\usepackage{geometry}
\usepackage{graphicx}
\usepackage{amsmath, amssymb}
\geometry{margin=1in}
\begin{document}
\begin{center}
Astronomy - Graduate Level Assessment 14 \
Total Marks: 40 \
Answer all questions.
\end{center}

\section*{Multiple Choice (10 marks)}
\begin{enumerate}[label=\arabic*.]
\item The belts and zones of Jupiter are
\begin{enumerate}[label=(\alph*)]
    \item names for the layers of gaseous and metallic hydrogen deep within the planet.
    \item alternating bands of rising and falling air at different latitudes.
    \item alternating regions of charged particles in Jupiter's magnetic field.
    \item regions of the plasma torus created by ions from Io's volcanoes
\end{enumerate}
\item When traveling north from the United States into Canada you'll see the North Star (Polaris) getting \_\_\_\_\_\_\_\_\_.
\begin{enumerate}[label=(\alph*)]
    \item Brighter
    \item Dimmer
    \item Higher in the sky
    \item Lower in the sky
\end{enumerate}
\item Suppose the angular separation of two stars is smaller than the angular resolution of your eyes. How will the stars appear to your eyes?
\begin{enumerate}[label=(\alph*)]
    \item You will not be able to see these two stars at all.
    \item You will see two distinct stars.
    \item The two stars will look like a single point of light.
    \item The two stars will appear to be touching looking rather like a small dumbbell.
\end{enumerate}
\item Where are the Trojan asteroids located?
\begin{enumerate}[label=(\alph*)]
    \item in the center of the asteroid belt
    \item on orbits that cross Earth's orbit
    \item surrounding Jupiter
    \item along Jupiter's orbit 60^\ensuremath{\circ} ahead of and behind Jupiter
\end{enumerate}
\item Kirkwood gaps are observed in the main asteroid belt including at the position(s) where:
\begin{enumerate}[label=(\alph*)]
    \item asteroids would orbit with a period half that of Jupiter
    \item asteroids would orbit with a period twice that of Jupiter
    \item asteroids would orbit with a period twice that of Mars
    \item A and B
\end{enumerate}
\end{enumerate}

\section*{True / False (10 marks)}
\textit{Answer True or False}
\begin{enumerate}[label=\arabic*.]
\setcounter{enumi}{5}
\item True or False: Mercury's atmosphere is primarily lost through volcanic heating, which continually replenishes it.
\item True or False: Organisms living deep in the oceans around seafloor volcanic vents and in hot springs most resemble the common ancestor of all life.
\item True or False: Moons significantly contribute to the stability of particles within rings, create gaps, and exert gravitational effects at ring edges.
\item True or False: Jupiter's magnetic field is approximately 20,000 times stronger than Earth's magnetic field.
\item True or False: A solar day on Planet X is approximately 100 Earth days due to its prograde rotation and orbital synchronization.
\end{enumerate}

\section*{Long Form Response (20 marks)}
\textit{Answer each question with a clear and complete explanation.}
\begin{enumerate}[label=\arabic*.]
\setcounter{enumi}{10}
\item Discuss the factors influencing the acceleration of a truck on Mars compared to Earth, assuming no friction. What conclusions can be drawn about the role of gravitational force in this scenario?
\item Discuss the concept of black bodies in astronomy, including their definition, significance in calculating stellar radiation, and implications for understanding stellar properties.
\end{enumerate}

\vfill
\noindent\textit{End of Paper}
\end{document}